% Copyright (C) 2019-2026, Lux Industries Inc. All rights reserved.
% pulsard: the offline threshold ML-DSA (FIPS-204) signer for the Pulsar lane.
% Design of the dealerless TALUS path, its fail-closed boundary, and the
% interop+trust contract a concrete engine must satisfy.
\documentclass[11pt]{article}

\usepackage{amsmath,amssymb,amsthm}
\usepackage{enumitem}
\usepackage[margin=1in]{geometry}
\usepackage{hyperref}
\usepackage{xcolor}

\newcommand{\HighBits}{\mathsf{HighBits}}
\newcommand{\LowBits}{\mathsf{LowBits}}
\newcommand{\MakeHint}{\mathsf{MakeHint}}
\newcommand{\FindHint}{\mathsf{FindHint}}
\newcommand{\Verify}{\mathsf{Verify}}
\newcommand{\Sign}{\mathsf{Sign}}
\newcommand{\ctx}{\mathit{ctx}}
\newcommand{\subj}{\mathit{subject}}
\newcommand{\Ring}{R_q}
\newcommand{\norm}[1]{\lVert #1 \rVert_\infty}

\theoremstyle{definition}
\newtheorem{contract}{Contract}
\theoremstyle{plain}
\newtheorem{claim}{Claim}

\title{\texttt{pulsard}: Offline Threshold ML-DSA for the Lux Pulsar Lane\\
\large Design of the dealerless TALUS path and its fail-closed boundary}
\author{Lux Industries --- Cryptography}
\date{}

\begin{document}
\maketitle

\begin{abstract}
\texttt{pulsard} is the offline signer that produces the Lux Quasar
\emph{Pulsar} finality lane: \emph{one} ordinary FIPS-204 ML-DSA-65 signature
over a $32$-byte subject, verifiable on chain by \texttt{warp.VerifyPulsar} (a
plain \(\Verify\) under a group public key). This document specifies the
intended \emph{dealerless} threshold construction (TALUS), states precisely
where the implementation is deliberately \emph{fail-closed}, and fixes the
interop and trust contract that any concrete threshold engine must satisfy
before it may be wired behind the \texttt{ThresholdEngine} SPI. It is written to
be adversarially reviewed: the honest position is that the dealerless threshold
crypto is research-grade and is therefore \emph{not} shipped as working code.
\end{abstract}

\section{Scope and the produce/verify asymmetry}

The chain's trust anchor is a single, standard object:
\[
\Verify\bigl(\mathsf{pk}_{\text{group}},\ \subj,\ \ctx,\ \sigma\bigr)
\quad\text{with}\quad
\ctx = \texttt{"LUX-QUASAR-PULSAR-MLDSA65-v1"},
\]
where \(\mathsf{pk}_{\text{group}}\) is an ordinary ML-DSA-65 public key and
\(\sigma\) is an ordinary FIPS-204 signature. Everything below the verification
boundary --- nonce DKG, commitment carry, carry elimination, secure comparison,
blame, aggregation --- exists only to \emph{produce} \(\sigma\) from key shares
without ever reconstructing the secret. The verifier never learns the protocol
that produced \(\sigma\); it checks a single ML-DSA signature. \texttt{pulsard}
realises the \emph{produce} side; \texttt{warp} realises \emph{verify}; the two
share exactly one definition of the key-era record (\texttt{warp.PulsarKeyEra}).

\paragraph{The release gate.} Independently of which engine produces \(\sigma\),
\texttt{pulsard} re-runs \texttt{warp.VerifyPulsar} on the assembled evidence
before returning it (\texttt{Signer.ThresholdSign} $\to$ \texttt{ReleaseGate}).
A signature that does not verify under the chain's own verifier is discarded.
Thus correctness of \emph{emitted} evidence is a property of \texttt{pulsard},
not of the engine: a buggy or malicious engine can cause liveness failure, never
the emission of invalid evidence.

\section{Why threshold ML-DSA is hard}

ML-DSA (FIPS-204, Dilithium lineage) signs by sampling a masking nonce
\(y\), computing the commitment \(w = A y\), deriving the challenge
\(c = H(\mu \,\|\, w_1)\) where \(w_1 = \HighBits(w)\), and releasing
\(z = y + c\,s_1\) together with a hint \(h\). The map is \emph{not} linear in
the nonce: \(\HighBits\) and the low-order residual \(r_0=\LowBits(w)\) are
rounding operations. Consequently there is no homomorphic nonce commitment, and
the FROST blueprint --- additively share the nonce, sum partial signatures ---
\emph{does not} yield a valid ML-DSA signature. This is the wall every
threshold-ML-DSA attempt hits.

\section{The TALUS construction (intended dealerless path)}
\label{sec:talus}

TALUS (Kao, \emph{Threshold ML-DSA with One-Round Online Signing via Boundary
Clearance and Carry Elimination}, arXiv:2603.22109) goes \emph{around} the wall
rather than through it. \texttt{pulsard} models its phases as a control-plane
state machine (\texttt{protocol.go}, \texttt{session.go}); the cryptographic
content of each phase is the engine's responsibility.

\begin{description}[leftmargin=1.6em]
  \item[Shamir nonce DKG (\textsc{PhaseNonceDKG}).] A dealerless, one-time joint
  nonce \(\bar y = \sum_i \lambda_i y_i\) is Shamir-shared so that \emph{no}
  node ever forms \(\bar y\). (\texttt{NonceDKGDeal}.)

  \item[Boundary Clearance Condition (\textsc{PhaseBCC}).] Accept only nonces
  whose commitment keeps \(\norm{r_0} < \gamma_2 - \beta\), so the secret shift
  \(c\,s_2\) cannot cross a \(\HighBits\) boundary and the hint becomes
  \emph{publicly} computable from \(w' = A z - c\,t_1 2^d\) via \(\FindHint\).
  The acceptance probability is \(\approx 31.7\%\); a non-clear nonce is
  discarded. (\texttt{BCCShare}.)

  \item[Carry Elimination Framework (\textsc{PhaseCEF}).] Compute
  \(w_1 = \HighBits(A\bar y)\) over the secret-shared nonce \emph{without}
  reconstructing \(\bar y\) and \emph{without} revealing \(w_0=\LowBits(w)\)
  (revealing \(w_0\) leaks the key residual \(w'-w = c\,t_0 - c\,s_2\); this is
  the \textsc{W-LEAK} hazard). (\texttt{CEFCarryShare}.)

  \item[Secure-comparison sub-protocol (\textsc{PhaseCSCP}).] Realise the
  \(\HighBits\) boundary count as a secure comparison over secret-shared \(w\),
  so that no node or process forms \(w_0\) or \(w\) even transiently. This is the
  step that closes the \textsc{W-LEAK} residual. Information-theoretic privacy of
  the degree-\(2(T-1)\) products formed by the BGW multiplication substrate
  requires, for reconstruction threshold \(T \ge 3\), a committee of
  \(N \ge 2T-1\) (TALUS Thm.~10.1; \texttt{Committee.MinPartiesMPC}).
  (\texttt{CSCPShare}.)

  \item[Blame (\textsc{PhaseBlame}).] Identifiable abort: a detected fault is
  attributed to a specific node before the session aborts. (\texttt{BlameAccusation}.)

  \item[Aggregate (\textsc{PhaseAggregate}).] The single \emph{online} broadcast
  round: derive \(c = H(\mu \,\|\, w_1)\); each signer broadcasts
  \(z_i = \lambda_i y_i + c\,\lambda_i s_{1,i}\); the coordinator sums
  \(z = \sum_i z_i\), recovers \(h\) from public data, and assembles a stock
  FIPS-204 signature \((\tilde c, z, h)\). (\texttt{PartialZ}.)
\end{description}

\paragraph{Two trust profiles.} The online path and the emitted signature are
byte-identical across both; only the custody of \(w\) during preprocessing
differs. \textsc{Talus-TEE} lets a trusted coordinator/TEE form \(w\) and the
BCC filter directly (no honest-majority bound). \textsc{Talus-MPC} is TEE-free:
\(w_1\) is computed by CEF/CSCP over shares (honest majority, \(N\ge 2T-1\) for
\(T\ge 3\)). \emph{Only \textsc{Talus-MPC} delivers the permissionless property
the Pulsar lane advertises.}

\section{Security argument and its boundary}

\paragraph{What the signature inherits.} The output is a standard FIPS-204
signature. Its unforgeability is exactly ML-DSA EUF-CMA, which reduces to
Module-LWE / Module-SIS in the (Q)ROM. BCC and CEF change only \emph{how} the
same \((\tilde c, z, h)\) is produced, never \emph{what} a verifier checks.

\paragraph{What the protocol must additionally guarantee} (and what is
\emph{not} yet established in shipped code):
\begin{enumerate}[leftmargin=1.6em]
  \item \textbf{Zero key leakage.} The transcript must not reveal \(s_1, s_2,
  t_0\) or \(w_0\). The \textsc{W-LEAK} hazard is real; the CSCP is the
  mitigation. A semi-honest-only CSCP does not meet the malicious threat model.
  \item \textbf{Dealerless keygen.} Permissionless safety requires the group key
  itself to come from a dealerless DKG. A trusted-dealer keygen reintroduces a
  single point of compromise and is \emph{not} the advertised property.
  \item \textbf{Identifiable abort under honest majority.} A malicious minority
  must not be able to halt progress without attribution.
  \item \textbf{Transcript indistinguishability (caveat).} The threshold signing
  transcript may be \emph{distinguishable} in distribution from a single-party
  ML-DSA transcript (masked CEF broadcasts and per-party \(z_i\) are extra
  observables), even though the \emph{final signature} is byte-standard. This is
  a known caveat to state, bound, and review --- not to paper over.
\end{enumerate}

\section{The fail-closed boundary}

\texttt{pulsard} ships the control plane (state machine, session/nonce plumbing,
key-era records, release gate) as working, tested code, and the threshold
\emph{crypto} as a fail-closed stub:
\begin{itemize}[leftmargin=1.6em]
  \item The default \texttt{ThresholdEngine} is \texttt{Unimplemented()}; every
  crypto operation returns \texttt{ErrThresholdMLDSAUnimplemented}.
  \item No partial, hand-wavy threshold arithmetic is present. An honest
  fail-closed stub is correct; unsound ``looks-done'' threshold logic is a
  forgery surface and is worse than no code.
  \item The single-party \texttt{ReferenceDealer} is the only engine that
  produces a signature. It is a labelled dev/test \emph{footgun}: it holds the
  whole group secret and has \emph{no} threshold property. It exists to prove the
  verify integration and nothing else.
\end{itemize}

\section{The engine contract (review gate)}

A concrete dealerless engine (the natural candidate is the
\texttt{github.com/luxfi/pulsar} reference, which already emits CIRCL-compatible
FIPS-204 signatures) may be registered behind the SPI \emph{only} after all of
the following are established and reviewed.

\begin{contract}[Interop]
The engine's output $\sigma$ satisfies
\(\Verify(\mathsf{pk}_{\text{group}}, \subj, \ctx, \sigma)=\mathsf{true}\) under
\texttt{github.com/luxfi/crypto/pq/mldsa/mldsa65} with
\(\ctx=\texttt{LaneContext}\), for the \emph{same} resolved
\(\mathsf{pk}_{\text{group}}\) bytes the chain resolves. This is mechanically
enforced by \texttt{ReleaseGate}; it must also be proven for the engine's own
key encoding (CIRCL packed public key) so that liveness, not just safety, holds.
\end{contract}

\begin{contract}[Dealerless keygen]
\(\mathsf{pk}_{\text{group}}\) and the shares are produced by a dealerless DKG
(e.g.\ Pedersen/Gennaro over the relevant ring), with no party ever holding the
master secret. A trusted-dealer keygen disqualifies the engine from the
permissionless profile.
\end{contract}

\begin{contract}[Malicious security]
The CSCP/CEF realises secure \(\HighBits\) with malicious security and
identifiable abort under an honest majority (\(N\ge 2T-1\) for \(T\ge 3\)); the
transcript opens only \(\{\text{validity flags},\ \text{uniform masks},\ w_1\}\)
and provably no function of \(\{s_1,s_2,t_0,w_0\}\).
\end{contract}

\begin{contract}[Dual-PQ key independence]
When the Pulsar (ML-DSA) lane is combined with the Corona (Module-LWE) lane in
the dual-PQ AND-mode cert, the two group keys are generated from independent
randomness (no shared seed); otherwise dual-PQ collapses to single-PQ.
\end{contract}

\begin{contract}[Distribution caveat bounded]
The transcript-indistinguishability caveat of \S4 is either closed or its
adversarial advantage is bounded and accepted in the threat model.
\end{contract}

Until Contracts~1--5 hold, the dealerless engine stays unregistered and
\texttt{pulsard} stays fail-closed. This is the deliberate, honest state of the
art: the service, the verify integration, and the protocol structure are real;
the dealerless threshold crypto is pending review.

\section*{References}
\begin{itemize}[leftmargin=1.6em]
  \item NIST FIPS 204, \emph{Module-Lattice-Based Digital Signature Standard}, 2024.
  \item M.~Kao, \emph{TALUS: Threshold ML-DSA with One-Round Online Signing via
  Boundary Clearance and Carry Elimination}, arXiv:2603.22109.
  \item Ben-Or, Goldwasser, Wigderson, \emph{Completeness Theorems for
  Non-Cryptographic Fault-Tolerant Distributed Computation}, STOC 1988 (BGW).
  \item Gennaro, Jarecki, Krawczyk, Rabin, \emph{Secure Distributed Key
  Generation for Discrete-Log Based Cryptosystems}, EUROCRYPT 1999.
  \item \texttt{github.com/luxfi/pulsar} --- the TALUS reference implementation
  (BCC/CEF/CSCP), and its \texttt{BLOCKERS.md}/\texttt{AUDIT-2026-06.md}.
\end{itemize}

\end{document}
